\begin{lemma}[Perfect base extensions yield perfect sub-super-sequences]
Let $f:F\to Q$ be a super-sequence on $X$, and let $h$ be a base
multi-sequence satisfying the value equation for the front extension.  If
$z\in\mathsf{BaseIncSeq}(X)$ and the global-coordinate restriction
$x\mapsto h(z\circ x)$ is perfect, then $f$ has a perfect restriction to a
sub-front on the exact subbase $Y=\operatorname{range}(z)$.
\end{lemma}
\begin{proof}
Write $Y=\operatorname{range}(z)$.  By the definition of
\noderef{BaseIncSeq}, $z$ is an element of \mathsf{IncSeq} and
$Y\subseteq X$.  The set $Y$ is infinite: if it had $N$ elements, then the
$N+1$ values $z(0),\ldots,z(N)$ would be pairwise distinct by the order
embedding property in \noderef{IncSeq}, contradicting that they all belong
to a set with only $N$ elements.  Thus $Y$ is an infinite subset of $X$.

Define
\[
 F'=F\mathbin{|}Y=\mathsf{FrontRestrict}(F,Y).
\]
The definition of \noderef{FrontRestrict} gives the inclusion
$F'\subseteq F$: if $s\in F'$, then $s\in F$ and every member of $s$ lies
in $Y$.  Apply \noderef{FrontRestriction} to the front
$f.\mathsf{front}$ recorded by \noderef{SuperSequence}, the inclusion
$Y\subseteq X$, and the infinitude just proved.  We obtain
\[
 hF':\mathsf{Front}(F',Y).
\]
We now prove the last component of the existential conclusion, namely
\[
 \mathsf{PerfectSuperSequence}(r,
   \mathsf{SuperSequenceRestrict}(f,F',Y,hF',F'\subseteq F)).
\]
Unfolding \noderef{PerfectSuperSequence}, fix $s,t\in F'$ with membership
proofs $h_s,h_t$ and assume
$\mathsf{FiniteShift}(s,t)$.  The definition \noderef{FrontRestrict}
supplies $s,t\in F$, and \noderef{FiniteShift} supplies an
$x\in\mathsf{IncSeq}$ together with
\[
 p_s:\mathsf{ProperPrefixSet}(s,\operatorname{range}(x)),\qquad
 p_t:\mathsf{ProperPrefixSet}(t,\operatorname{range}(\mathsf{ShiftMap}(x))).
 \tag{1}
\]

We first replace this witness by one whose range is contained in $Y$.  The
empty case must be separated because there need not be an element of $Y$
below an arbitrary first element of the original witness.  If exactly one
of $s,t$ were empty, the nonempty one would have the empty set as an
initial segment: take its least element in the definition of
\noderef{InitialSegment}.  This contradicts the prefix-free field of the
front \noderef{Front}, since both sets belong to $F$.  Hence either
$s=t=\emptyset$, which is also the only possibility for a trivial front,
or both are nonempty.

Suppose first that $s=t=\emptyset$.  By \noderef{InfiniteSetEnumeration}
choose $y\in\mathsf{IncSeq}$ with $\operatorname{range}(y)=Y$.  The first
value $y(0)$ is the least value in its range, and the first value of
$\mathsf{ShiftMap}(y)$ is $y(1)$ and is the least value in that shifted
range, by \noderef{IncSeq} and \noderef{ShiftMap}.  Therefore the two
empty sets satisfy
\[
 \mathsf{ProperPrefixSet}(\emptyset,\operatorname{range}(y)),\qquad
 \mathsf{ProperPrefixSet}(\emptyset,
       \operatorname{range}(\mathsf{ShiftMap}(y))).                \tag{2}
\]

It remains to treat the nonempty case.  Unpack the two assertions in (1)
using \noderef{ProperPrefixSet}.  There are cutoffs
$a\in\operatorname{range}(x)$ and
$b\in\operatorname{range}(\mathsf{ShiftMap}(x))$ such that
\[
 s=\{k\in\operatorname{range}(x):k<a\},\qquad
 t=\{k\in\operatorname{range}(\mathsf{ShiftMap}(x)):k<b\}.          \tag{3}
\]
Choose $i,j\in\mathbb N$ with $a=x(i)$ and
$b=\mathsf{ShiftMap}(x)(j)=x(j+1)$; the last equality is the definition
of the shift in \noderef{ShiftMap}.  Since $x$ is an order embedding, (3)
becomes
\[
 s=\{x(k):k<i\},\qquad t=\{x(k):1\leq k\leq j\}.                 \tag{4}
\]
Indeed, order reflection says $x(k)<x(i)$ exactly when $k<i$, and the
same argument applied to the shifted sequence gives the second equality.

Put $L=\max(i,j+1)$.  In the nonempty case $i\geq1$.  Every entry
$x(k)$ with $k<L$ lies in $s\cup t$: entries with $k<i$ lie in $s$, and
the remaining entries have $i\leq k\leq j$ and lie in $t$.  Since
$s,t\in F'$ and \noderef{FrontRestrict} says that members of $F'$ are
subsets of $Y$, the finite set
\[
 P=\{x(k):k<L\}
\]
is contained in $Y$.  Let $p=x(L-1)$.  The tail
\[
 Y_{>p}=\{n\in Y:p<n\}
\]
is infinite.  Otherwise $Y$ would be the union of the finite initial
segment $\{n:n\leq p\}$ and the finite set $Y_{>p}$, contradicting the
infinitude of $Y$.  Use \noderef{InfiniteSetEnumeration} to choose
$v\in\mathsf{IncSeq}$ with range $Y_{>p}$.  Define $y$ by concatenating
the finite list
\[
 x(0),x(1),\ldots,x(L-1)
\]
with the increasing sequence $v$.  Every entry of the first list lies in
$Y$, and every entry of $v$ is in $Y$ and is larger than $p$; consequently
the concatenation is an increasing embedding, hence an element of
\mathsf{IncSeq} by \noderef{IncSeq}, and
$\operatorname{range}(y)\subseteq Y$.

We check both cutoffs, including the boundary cases.  For $s$, if
$i<L$, then $i\leq j$ and $i\geq1$, so $x(i)\in t\subseteq Y$ and is
the $i$-th entry of $y.  If $i=L$, the first value of $v$ is instead the
cutoff.  In either case, the entries of $y$ strictly below the chosen
cutoff are exactly $x(0),\ldots,x(i-1)$, which is $s$ by (4).  Thus
\[
 \mathsf{ProperPrefixSet}(s,\operatorname{range}(y)).               \tag{5}
\]
For $t$, if $j+1<L$, then $j+1<i$ and $x(j+1)\in s\subseteq Y$ is an
entry of $y$; if $j+1=L$, the first value of $v$ is the cutoff.  After
shifting $y$, the entries strictly below that cutoff are exactly
$x(1),\ldots,x(j)$, namely $t$.  Hence
\[
 \mathsf{ProperPrefixSet}(t,
       \operatorname{range}(\mathsf{ShiftMap}(y))).                 \tag{6}
\]
Equations (5) and (6) give the required subbase-preserving witness in the
nonempty case, while (2) gives it in the empty case.

We now factor this $y$ through the fixed enumeration $z$.  Since
$\operatorname{range}(y)\subseteq Y=\operatorname{range}(z)$, for every
$n$ there is a unique $u(n)$ with $z(u(n))=y(n)$, uniqueness following
from injectivity in \noderef{IncSeq}.  If $n<m$, then
$y(n)<y(m)$, so order reflection for $z$ gives $u(n)<u(m)$.  Thus
$u\in\mathsf{IncSeq}$, and pointwise $y(n)=z(u(n))$ gives
\[
 y=\mathsf{IncSeqComp}(z,u).                                      \tag{7}
\]
The composition statement \noderef{IncSeqComp} also shows that
$\mathsf{IncSeqComp}(z,u)$ and
$\mathsf{IncSeqComp}(z,\mathsf{ShiftMap}(u))$ are increasing
embeddings.  Their ranges are contained in $\operatorname{range}(z)$,
and hence in $X$ by \noderef{BaseIncSeq}.  They therefore define
\[
 x_0=\langle\mathsf{IncSeqComp}(z,u),\ \operatorname{range}\subseteq X\rangle,
 \qquad
 x_1=\langle\mathsf{IncSeqComp}(z,\mathsf{ShiftMap}(u)),\ \operatorname{range}\subseteq X\rangle
 \tag{8}
\]
as elements of \mathsf{BaseIncSeq}(X).  By (7) and
\noderef{RestrictionShift},
\[
 x_0.1=y,\qquad
 x_1.1=\mathsf{ShiftMap}(y).                                      \tag{9}
\]

The first components of the two \noderef{FrontRestrict} membership proofs
give $s,t\in F$.  Combining these memberships with (5) and (6), or with
(2) in the empty case, and using \noderef{FrontPrefix}, we obtain
\[
 q_s:\mathsf{FrontPrefix}(F,s,x_0.1),\qquad
 q_t:\mathsf{FrontPrefix}(F,t,x_1.1).                              \tag{10}
\]
Apply the extension equation in the theorem statement first to
$(x_0,s,q_s)$ and then to $(x_1,t,q_t)$.  It gives the two dependent-value
equalities
\[
 h(x_0)=f.\mathsf{value}\langle s,q_s.1\rangle,qquad
 h(x_1)=f.\mathsf{value}\langle t,q_t.1\rangle.                  \tag{11}
\]
By the definition of \noderef{BaseRestrict}, the two left sides in (11)
are respectively
\[
 (\mathsf{BaseRestrict}(h,z))(u),\qquad
 (\mathsf{BaseRestrict}(h,z))(\mathsf{ShiftMap}(u)).                \tag{12}
\]
The hypotheses $q_s.1$ and $q_t.1$ are proofs that $s,t\in F$; proof
irrelevance identifies them with the membership proofs obtained from
$F'\subseteq F$.  Thus (11) also identifies the right sides with the
values of the restricted super-sequence at the dependent subtype points
$\langle s,h_s\rangle$ and $\langle t,h_t\rangle$.

Finally, unfold \noderef{PerfectMultiSequence} in the hypothesis $hp$ and
instantiate it at $u$.  Using (12), it yields
\[
 r\bigl(h(x_0),h(x_1)\bigr).                                      \tag{13}
\]
Unfolding \noderef{SuperSequenceRestrict} replaces the two values of the
restricted super-sequence by the two right sides of (11), and (11) then
rewrites (13) as
\[
 r\Bigl(\mathsf{SuperSequenceRestrict}(f,F',Y,hF',F'\subseteq F).value
       \langle s,h_s\rangle,
       \mathsf{SuperSequenceRestrict}(f,F',Y,hF',F'\subseteq F).value
       \langle t,h_t\rangle\Bigr).
\]
This is the required implication for arbitrary $s,t$ in $F'$, so the
unfolded \noderef{PerfectSuperSequence} predicate holds.  Together with
$hF'$ and the inclusion supplied by \noderef{FrontRestrict}, this is the
claimed existential conclusion.
\end{proof}
